\section{The physical pulse curve and the conditional crossing}\label{sec:pulse-crossing}
\subsection{A correlated wide-parameter pulse family}

Let
\[
 J_0=\frac{2409}{8000},\qquad
 h=\frac3{40000},\qquad J=J_0+h\zeta,\quad -1\le\zeta\le1.
\]
For the pulse solution \(z(t,J)\), propagate the scaled coefficients
\begin{equation}
 b_k(t)=\frac{h^k}{k!}\partial_J^kz(t,J_0),
 \qquad 0\le k\le4,                                    \label{eq:scaled-jet}
\end{equation}
as one correlated system.  Because the vector field is cubic, coefficient
products are exact finite convolutions; the degree-five through degree-twelve
terms form the forcing for the remainder
\begin{equation}
 R_5(t,\zeta)=z(t,J_0+h\zeta)-\sum_{k=0}^4b_k(t)\zeta^k.
                                                               \label{eq:R5}
\end{equation}
This scaled representation retains the cancellation that is destroyed by
independent derivative boxes.

The \(192\)-bit directed Taylor--Bernstein computation covers all \(1152\)
time cells on \(0\le t\le24\sqrt5\).  With time degree \(16\), every
coefficient guide and every remainder tube closes.  The certified maxima
are
\begin{align}
 \max\norm{\text{coefficient error}}_P
   &\le8.5725360731\times10^{-19},\\
 \max\norm{\text{degree }5\text{--}12\text{ forcing}}_P
   &\le2.0594765514\times10^{-9},\\
 \max\norm{R_5}_P
   &\le1.7206344728\times10^{-8},                       \label{eq:wide-jet-bounds}\\
 \max\abs{(R_5)_v}
   &\le3.5937291129\times10^{-9}.
\end{align}
The complete computation has been independently regenerated from its fixed
inputs.  Hence the wide fixed-time parameter family is proved.  These bounds
do not yet include a moving event time.

\subsection{Event alignment and the scalar crossing}

Let \(g(t,\xi,J)\) define a transverse event and suppose
\begin{equation}
 g(\tau(\xi,J),\xi,J)=0,\qquad
 \abs{\partial_tg}\ge\sigma>0.                         \label{eq:event-speed}
\end{equation}
The implicit-function theorem gives
\begin{equation}
 \partial_J\tau=-\frac{\partial_Jg}{\partial_tg}.       \label{eq:event-first-jet}
\end{equation}
Higher event jets follow recursively by differentiating
\(g(\tau(\xi,J),\xi,J)=0\).  If
\(z_t(\theta)=z(t+\theta)\), the common-event retained reduced history is
\begin{equation}
 \begin{split}
 B_\xi(J):=K_\xi(J):=\Pi_Yz_{\tau(\xi,J)}
   =\bigl(&\theta\longmapsto
       v(\tau(\xi,J)+\theta,\xi,J),\\
          &w(\tau(\xi,J),\xi,J)\bigr)\in Y,
 \qquad-r\le\theta\le0.                                \label{eq:event-history}
 \end{split}
\end{equation}
Thus \(B_\xi\) is the notation used in the completion theorem, while
\(K_\xi\) is retained below for compatibility with the certified
event-alignment chain; neither denotes the release history \(R_\xi\).
Its two reduced components have derivatives
\begin{align*}
 \partial_JK_{\xi,v}(J)(\theta)
 &=v_J(\tau+\theta,\xi,J)
   +\dot v(\tau+\theta,\xi,J)\partial_J\tau,
   &&-r\le\theta\le0,\\
 \partial_JK_{\xi,w}(J)
 &=w_J(\tau,\xi,J)+\dot w(\tau,\xi,J)\partial_J\tau.
\end{align*}
Thus event alignment is a history translation problem and cannot be
replaced by an endpoint correction.

For the physical pulse family in \cref{eq:scaled-jet}, let
\(h_C(\phi)=\phi_v(0)-V_i(0)\), where the exact phase-zero inner-orbit
voltage is source-bound by
\begin{equation}
 0.9053838432821200<V_i(0)<0.9054038432821201.          \label{eq:route-c-level}
\end{equation}
The wider Fourier candidate section used during orbit discovery is not used
as a theorem input.

\begin{theorem}[Uniform Route--C event alignment]\label{thm:route-c-event}
For every
\(J\in[6021/20000,753/2500]=[0.30105,0.30120]\), the exact pulse solution
has exactly one positive crossing of \(h_C=0\) in
\begin{equation}
 I_C=\left[\frac{555\sqrt5}{24},
               1+\frac{546\sqrt5}{24}\right].         \label{eq:route-c-bracket}
\end{equation}
On the whole parameter--time box its event speed satisfies
\begin{equation}
 0.2133519018734632<\dot v<0.2791999723595853.          \label{eq:route-c-speed}
\end{equation}
Writing \(J=J_0+h\zeta\), the event time has a fourth-order polynomial
\(\widehat T_4(\zeta)\) such that
\begin{equation}
 \abs{T(\zeta)-\widehat T_4(\zeta)}\le10^{-4},
 \qquad -1\le\zeta\le1.                              \label{eq:event-time-tube}
\end{equation}
At \(J_0\),
\begin{equation}
 51.79830610110210<T(0)<51.79838713426455.             \label{eq:center-event-time}
\end{equation}
The scaled coefficients
\(\widehat T_4(\zeta)=T_0+\sum_{k=1}^4a_k\zeta^k\) satisfy
\begin{align}
 a_1&\in[2.9068862680,2.9074011687]\times10^{-2},&
 a_2&\in[1.9583818528,1.9612556829]\times10^{-4},\notag\\
 a_3&\in[1.9208606177,1.9385338925]\times10^{-6},&
 a_4&\in[2.2095125644,2.2986280783]\times10^{-8}.       \label{eq:event-time-coefficients}
\end{align}
Moreover the translated histories
\(\mathcal K(\zeta)=(v(T(\zeta)+\cdot,J),w(T(\zeta),J))\in Y\) obey
\begin{equation}
 \norm{\mathcal K(\zeta)-B_c(\zeta)}_Y
 \le0.008199931248423,                                 \label{eq:event-history-tube}
\end{equation}
where the continuous fourth-order fixed-time reference family is
\[
 B_c(\zeta)=
 \left(\theta\longmapsto\sum_{k=0}^4b_{k,v}(T_0+\theta)\zeta^k,\,
                   \sum_{k=0}^4b_{k,w}(T_0)\zeta^k\right).
\]
\end{theorem}

The proof combines the source-bound fixed-time family with cellwise
Bernstein endpoint and vector-field bounds.  The gap is uniformly negative
at \(\widehat T_4-10^{-4}\), with upper bound
\(-6.4436\times10^{-7}\), and uniformly positive at
\(\widehat T_4+10^{-4}\), with lower bound
\(6.4342\times10^{-7}\).  The positive speed then gives existence and
uniqueness.  This proves neither that the event is the third post-release
crossing nor that \(\mathcal K(\zeta)\) lies on the inner stable sheet.  A
fourth-order \(Y\)-valued history jet is not claimed.  The first derivative,
including the moving-event contribution, is certified next.

\begin{theorem}[Event-aligned first pulse derivative]
\label{thm:event-first-derivative}
Put \(K(J)=\mathcal K((J-J_0)/h)\).  The map
\[
 K:[0.30105,0.30120]\longrightarrow Y
\]
is continuously differentiable.  If \(W=\partial_\zeta z\),
\(C=\partial_\zeta\sum_{k=0}^4b_k\zeta^k\), and \(E=W-C\), the directed
comparison for \(E\) closes on all \(1152\) cells with maximum weighted
radius
\begin{equation}
 8.7288930458\times10^{-8}.                            \label{eq:stage5d-radius}
\end{equation}
The event derivative satisfies
\begin{equation}
 T_J\in[336.6243028349,456.5740938122].                \label{eq:TJ-bound}
\end{equation}
On the complete event history,
\begin{align}
 D_JK_v(0)&=0,\notag\\
 D_JK_w&\in[-17.350656459,-10.142798861],\notag\\
 \norm{D_JK}_Y&\le142.200202974.                       \label{eq:DJK-bound}
\end{align}
For the fixed center Route--C functional from
\cref{eq:fourier-reversal,eq:fq-bound}, the presently directed consequence
is only
\begin{equation}
 \abs{f(D_JK)}\le1017.273043.                          \label{eq:DJK-action-modulus}
\end{equation}
No sign or exclusion of zero is asserted in
\cref{eq:DJK-action-modulus}.
\end{theorem}

The proof does not differentiate the fixed-time remainder bound.  The exact
error equation is
\[
 E'=DF(z)E+\{DF(z)-DF(B)\}C
       +\partial_\zeta\operatorname{Tail}_{\ge5}(B),
 \qquad B=\sum_{k=0}^4b_k\zeta^k.
\]
The state remainder enters only through the correlated linearization
difference.  A \(192\)-bit Taylor--Bernstein comparison on \(64\) parameter
shards closes every fixed-time cell.  The event identity
\[
 T_\zeta=-W_v(T,\zeta)/\dot v(T,\zeta)
\]
is then evaluated on \(128\) parameter shards, and the history chain rule
retains \(z_tT_J\) on the entire translated window.  Finally,
\cref{eq:DJK-action-modulus} uses the validated density total variation and
atom moduli.  This last operation deliberately loses orientation.

\begin{lemma}[Phase-invariant oriented quotient]
\label{lem:phase-invariant-quotient}
Let \(\ell\) be a nonzero complex representative of the left eigenspace of
the simple real unstable multiplier, and let \(q,y\in Y\) be real with
\(\ell(q)\ne0\).  Suppose one common directed computation gives
\[
 \abs{\ell(y)-a_0}\le r_a,
 \qquad \abs{\ell(q)-b_0}\le r_b<\abs{b_0}.
\]
With \(c_0=a_0/b_0\), the normalized action is real and obeys
\begin{equation}
 \frac{\ell(y)}{\ell(q)}
 \in[\Re c_0-r_c,\Re c_0+r_c],
 \qquad
 r_c=\frac{r_a+\abs{c_0}r_b}{\abs{b_0}-r_b}.          \label{eq:oriented-quotient}
\end{equation}
Equivalently, for any real \(c\), if the same row gives
\(\abs{\ell(y-cq)}\le r\) and \(\abs{\ell(q)}\ge b_->0\), then
\begin{equation}
 \frac{\ell(y)}{\ell(q)}
 \in[c-r/b_-,c+r/b_-].                                \label{eq:same-row-residual-quotient}
\end{equation}
\end{lemma}

Indeed the quotient is invariant under every nonzero complex rescaling of
\(\ell\).  Reality and simplicity imply that conjugation changes \(\ell\)
only by such a rescaling, so its action ratio on real histories is real.
The disk estimate in \cref{eq:oriented-quotient} follows by writing the
difference from \(a_0/b_0\) over the common denominator.
\Cref{eq:same-row-residual-quotient} follows from
\[
 \frac{\ell(y)}{\ell(q)}-c
 =\frac{\ell(y-cq)}{\ell(q)}.
\]
\Cref{eq:same-row-residual-quotient} is the preferred correlated form: choose
a source-bound center \(c\), form \(D_JK-cq\) as one complete-history object,
apply the same Fourier row and tail enclosure, and divide only after the
correlated residual action has been bounded.  This is the first-variation
analogue of the successful same-row deflation in
\cref{prop:base-uu}.

\begin{lemma}[Physical right-gauge correction]
\label{lem:right-gauge}
Let \(\widetilde q\) be a complex Grushin-normalized right eigencolumn for a
simple real multiplier of a real return operator.  Choose a real evaluation
\(\chi\) with \(\chi(\widetilde q)\ne0\), and orient the physical unstable
history by
\[
 \gamma=\frac{\chi(\widetilde q)}
               {\abs{\chi(\widetilde q)}}\in\mathbb S^1,
 \qquad
 q=\gamma^{-1}\widetilde q .
\]
Then \(q\) is real and \(\chi(q)>0\).  For any raw left representative
\(\ell\) with \(\ell(\widetilde q)\ne0\), the real normalized unstable
functional and stable projection on real histories are
\begin{equation}
 f(y)=\gamma\frac{\ell(y)}{\ell(\widetilde q)},
 \qquad
 qf(y)=\widetilde q\frac{\ell(y)}{\ell(\widetilde q)}. \label{eq:right-gauge}
\end{equation}
In particular, the second expression is independent of both complex
left- and right-eigenvector phase gauges.
\end{lemma}

To prove the lemma, reality and simplicity imply
\(\widetilde q=\eta q_{\mathbb R}\) for a real eigenhistory
\(q_{\mathbb R}\) and some \(\eta\ne0\).  Orienting
\(q_{\mathbb R}\) by \(\chi(q_{\mathbb R})>0\) shows that
\(\gamma=\eta/\abs{\eta}\), so \(\gamma^{-1}\widetilde q\) is real.
The identities in \cref{eq:right-gauge} then follow by normalization.

\Cref{lem:right-gauge} is not optional in the executable certificate.  The
stored Grushin eigencolumn carries a nontrivial complex phase even though
the exact multiplier is real.  Dividing \(\ell(D_JK)\) directly by
\(\ell(\widetilde q)\) therefore produces a complex, right-gauge-dependent
number; it is not the oriented action in
\cref{lem:phase-invariant-quotient}.  The correlated residual must first use
the physical real history \(q=\gamma^{-1}\widetilde q\), or equivalently
insert the factor \(\gamma\) in \cref{eq:right-gauge}.  This correction moves
the diagnostic center from the recovery-atom shard near \(-14\) to the
normalized physical value near \(-252\).

\begin{theorem}[Fixed-center physical-phase action of the event derivative]
\label{thm:oriented-pulse-action}
Let \(q_{\rm phys}\) and \(f_{\rm phys}\) be the fixed center Grushin pair
oriented by \cref{lem:right-gauge}, with
\(f_{\rm phys}(q_{\rm phys})=1\).  For every
\(J\in[0.30105,0.30120]\),
\begin{equation}
 f_{\rm phys}(D_JK(J))\in
 [-258.746521015805,-245.253478984195]\subset(-\infty,0).
 \label{eq:oriented-pulse-action}
\end{equation}
The enclosure is uniform on the full pulse interval and contains the
event-time translation term.  It is an action of the fixed normalized
Route--C functional, not yet a stable-sheet-gap derivative.
\end{theorem}

The source-bound calculation takes \(c_*=-252\), forms the complete history
\(Y_*(J)=D_JK(J)-c_*q_{\rm phys}\), and keeps the recovery atom and both
delayed-density pieces in one directed action.  The current voltage atom
vanishes exactly because \(D_JK_v(0)=(q_{\rm phys})_v(0)=0\).  On \(128\)
pulse-parameter shards and \(512\) exact history cells, the two additive
joint acceptance envelopes give
\begin{align}
 r_{\rm guide}&\le0.002869679658061285,\notag\\
 r_{\rm measure}&\le1.899608044833\times10^{-6},\notag\\
 r_A&\le0.002871579266106117.                         \label{eq:stage5e-budget}
\end{align}
The same-row denominator satisfies
\begin{equation}
 |\ell(q_{\rm phys})|\ge0.0004256385267872229.       \label{eq:stage5e-denominator}
\end{equation}
Thus \cref{eq:same-row-residual-quotient} gives radius at most
\(6.746521015805\) about \(c_*\), proving
\cref{eq:oriented-pulse-action}.  All arithmetic is outward at \(192\) bits;
the finite row, Neumann tail, orbit covariance, density-basis transfer,
convolution rounding, pulse remainder, event graph, and seams are covered
before the final division.

The executable biological-control contract imports exactly six outputs from
\cref{thm:route-c-event}: the section level, pulse interval, unique declared-
bracket event, speed interval, time-graph remainder, and common-event
\(Y\)-tube.  Its fields for \(J_c\), stable-gap signs, physical onset,
routing, and capture remain null or false; its optional interval-Newton
enclosure is likewise absent.  Event alignment therefore becomes a proved
input to the control chain without silently promoting its next implication.

Once the graph \(\psi_\xi\) and normalized functional \(f_\xi\) are
available, let \(\phi_{i,\xi}^{\Sigma}\) be the inner-orbit base history on
the phase section and define the signed stable-sheet gap
\begin{equation}
 H(\xi,J)
 =f_\xi(K_\xi(J)-\phi_{i,\xi}^{\Sigma})
  -\psi_\xi\!\left[
   (\Id-q_\xi f_\xi)
   (K_\xi(J)-\phi_{i,\xi}^{\Sigma})\right].           \label{eq:stable-gap}
\end{equation}
For each fixed \(\xi\), the pair \((q_\xi,f_\xi)\) in this formula is held
fixed as \(J\) varies.
If \([F_J]\) encloses the oriented quotient
\(f_\xi(D_JK_\xi)\), \(\norm{D\psi_\xi}\le L_\psi\), and
\(\norm{(\Id-q_\xi f_\xi)D_JK_\xi}\le M_s\), then the exact correlated
derivative satisfies
\begin{equation}
\partial_JH(\xi,J)
\in[F_J]+[-L_\psi M_s,L_\psi M_s].                  \label{eq:gap-slope-budget}
\end{equation}
The oriented action and graph correction enter this estimate separately; the
norm in \cref{eq:DJK-bound} alone gives no slope.

At the center parameter, fix the exact reduced history space and Route--C
section
\[
 Y=C([-\tau_{\max},0],\mathbb R)\times\mathbb R,
 \qquad \Sigma=\{y\in Y:y_v(0)=0\},
\]
with the inherited max norm.  Let \(q_{\rm phys}\) and
\(f_{\rm phys}\) be the same fixed physically oriented Grushin pair used in
\cref{thm:oriented-pulse-action}, put
\(P_s=\Id-q_{\rm phys}f_{\rm phys}\), and set
\(\kappa(J)=K(J)-X_*\), where \(X_*\) is the inner-orbit reference history
on this section.  Thus the center specialization of \cref{eq:stable-gap} is
\(H(J)=f_{\rm phys}(\kappa(J))-\psi(P_s\kappa(J))\).

\begin{lemma}[Global affine splitting of the Route--C section]
\label{lem:route-c-affine-splitting}
Let \(q\in\Sigma\) and \(f\in\Sigma^*\) satisfy \(f(q)=1\), and put
\(P_s=\Id-qf\).  Then \(\operatorname{ran}P_s=\ker f\), and
\begin{equation}
 \mathcal C\colon\Sigma\longrightarrow\ker f\times\R,
 \qquad
 \mathcal C(y)=(P_sy,f(y))                              \label{eq:route-c-affine-chart}
\end{equation}
is a bounded linear isomorphism with inverse
\((z,u)\mapsto z+qu\).  Hence subtraction of \(X_*\) gives global
coordinates on the affine section \(X_*+\Sigma\).  If
\(\psi\colon D_s\subset\ker f\to\R\) represents the local stable sheet as
\begin{equation}
 X_*+\{z+q\psi(z):z\in D_s\},                          \label{eq:centered-stable-sheet}
\end{equation}
then, for every section history \(K\) with
\(P_s(K-X_*)\in D_s\),
\begin{equation}
 K\text{ belongs to this graph}
 \quad\Longleftrightarrow\quad
 f(K-X_*)-\psi(P_s(K-X_*))=0.                          \label{eq:affine-gap-equivalence}
\end{equation}
In particular, no additional numerical containment condition for an
ambient Route--C chart is required; only containment of the stable
coordinate in \(D_s\) is needed.
\end{lemma}

Indeed, \(f(P_sy)=0\) and \(y=P_sy+qf(y)\).  These identities prove both the
range statement and the two inverse relations in \cref{eq:route-c-affine-chart};
\cref{eq:affine-gap-equivalence} then follows from uniqueness of the split
coordinates.  This global affine coordinate statement does not enlarge the
local domain \(D_s\) on which the stable graph has been constructed.

\begin{theorem}[Centered endpoint coordinates, stable cone, and conditional slope]
\label{thm:conditional-stable-gap-slope}
On \(I_J=[J_-,J_+]=[0.30105,0.30120]\), source-bound complete-history
endpoint calculations give
\begin{equation}
 \begin{aligned}
 f_{\rm phys}(\kappa(J_-))&\in
   F_-:=[0.019677187,0.023055896],\\
 f_{\rm phys}(\kappa(J_+))&\in
   F_+:=[-0.018164491,-0.014783669],\\
 \norm{P_s\kappa(J_-)}_Y&\le0.008935972,
 &\norm{P_s\kappa(J_+)}_Y&\le0.008927665.
 \end{aligned}                                        \label{eq:stage5ga-endpoints}
\end{equation}
The direct norm of the same physical unstable column and the correlated
projection identity give
\begin{equation}
 \norm{q_{\rm phys}}_Y\le0.086274581,
 \qquad
 \sup_{J\in I_J}\norm{P_sD_J\kappa(J)}_Y\le5.979324. \label{eq:stage5gb-projection}
\end{equation}
A two-ended Banach-space cone, using both bounds in
\cref{eq:stage5ga-endpoints}, then proves
\begin{equation}
 P_s\kappa(I_J)
 \subset \overline B_Y\!\left(0,\frac{47}{5000}\right)
             \cap\ker f_{\rm phys}.                   \label{eq:stage5gb-cone}
\end{equation}
\Cref{prop:inner-terminal-stable-row} supplies the discrete stable-power input
\(K_s=1,\rho_s=0.1\) for a future graph transform.  It supplies neither the
nonlinear return tube nor any of the six uniform Hessian blocks.
If a quantitative centered \(C^1\) graph
\(\psi\colon D_s\subset\ker f_{\rm phys}\to\R\) is
constructed using this same \(Y,\Sigma,q_{\rm phys},f_{\rm phys},P_s\), if
\begin{equation}
 \overline B_Y\!\left(0,\frac{47}{5000}\right)
      \cap\ker f_{\rm phys}\subset D_s,
 \qquad
 \sup_{J\in I_J}\norm{D\psi(P_s\kappa(J))}\le16,     \label{eq:stage5gb-gate}
\end{equation}
then
\begin{equation}
 H'(J)\in[-354.415700,-149.584300]\subset(-\infty,0)
 \qquad(J\in I_J).                                    \label{eq:stage5gb-slope}
\end{equation}
The corresponding maximum admissible graph-derivative norm in this budget is
strictly larger than \(41.016926\).
\end{theorem}

For completeness, put \(z(J)=P_s\kappa(J)\) and
\(W=J_+-J_-=3/20000\).  From \cref{eq:stage5gb-projection}, the fundamental
theorem of calculus gives
\[
 \norm{z(J)}_Y\le
 \min\{E_-+5.979324(J-J_-),
        E_++5.979324(J_+-J)\},
\]
where \(E_-=0.008935972\) and \(E_+=0.008927665\) are the directed endpoint
bounds above.  The exact-rational
intersection of these two affine cones lies strictly inside \(I_J\), and its
directed maximum is at most
\(0.009380268<47/5000\), proving \cref{eq:stage5gb-cone}.  Combining
\cref{eq:gap-slope-budget,eq:stage5gb-projection} with
\cref{eq:oriented-pulse-action} proves the conditional interval
\cref{eq:stage5gb-slope}.  Thus the endpoint functional signs, the
full-interval stable-coordinate ball, and this conditional slope arithmetic
are proved.  The theorem does not supply the graph, either condition in
\cref{eq:stage5gb-gate}, endpoint graph heights, graph-adjusted stable-gap
signs, a selected crossing, onset, routing, or capture.

\begin{proposition}[Completion of the selected center-parameter crossing]
\label{prop:selected-center-crossing}
Let \(I_J=[J_-,J_+]=[0.30105,0.30120]\), and let
\(K\in C^1(I_J,X_*+\Sigma)\) be the single selected Route--C event branch enclosed in
\cref{thm:route-c-event,thm:event-first-derivative}.  Suppose a quantitative
centered \(C^1\) stable graph
\(\psi\colon D_s\subset\ker f_{\rm phys}\to\R\) is constructed in the identical
\(Y,\Sigma,q_{\rm phys},f_{\rm phys},P_s\) normalization and
\begin{equation}
 \overline B_Y\!\left(0,\frac{47}{5000}\right)
       \cap\ker f_{\rm phys}\subset D_s,
 \qquad
 \sup_{J\in I_J}\norm{D\psi(P_s\kappa(J))}\le16.     \label{eq:center-crossing-graph-data}
\end{equation}
Let \(F_-,F_+\) be the proved functional intervals in
\cref{eq:stage5ga-endpoints}.  Assume graph-height bounds
\(\eta_-,\eta_+\ge0\) such that
\begin{equation}
 \begin{aligned}
 \abs{\psi(P_s\kappa(J_-))}&\le\eta_-,\\
 \abs{\psi(P_s\kappa(J_+))}&\le\eta_+,
 \end{aligned}                                        \label{eq:center-crossing-endpoint-data}
\end{equation}
and that the graph-adjusted margins are strict:
\begin{equation}
 \inf F_- -\eta_->0,
 \qquad
 \sup F_+ +\eta_+<0.                                  \label{eq:center-crossing-margins}
\end{equation}
Then there is a unique
\begin{equation}
 J_c^{\rm sel}\in(J_-,J_+)
 \quad\text{such that}\quad
 H(J_c^{\rm sel})=0,                                  \label{eq:selected-center-root}
\end{equation}
and \(K(J_c^{\rm sel})\) lies on the local stable sheet in
\cref{eq:centered-stable-sheet}.  Moreover the intersection is transverse,
with \(H'(J_c^{\rm sel})\le-149.584300<0\).
\end{proposition}

To prove the proposition,
\cref{eq:stage5ga-endpoints,eq:center-crossing-endpoint-data,eq:center-crossing-margins}
give
\(H(J_-)>0>H(J_+)\).  The gap is continuous,
so the intermediate value theorem gives a zero.  Conditions
\cref{eq:center-crossing-graph-data} and
\cref{thm:conditional-stable-gap-slope} make \(H\) strictly decreasing on
all of \(I_J\), proving uniqueness and transversality; the stable-sheet
statement follows from the affine splitting lemma above.  Interval Newton
may subsequently replace \(I_J\) by a smaller certified enclosure of
\(J_c^{\rm sel}\), but it is not needed for either existence or uniqueness.

The conclusion of the proposition is one scalar at the center parameter
\(\xi_0=(1/4,1/200)\); it does not define a \(C^1\) function
\(J_c(\xi)\).  Such a function requires a common parameter neighborhood on
which \(K_\xi(J),X_{*,\xi},q_\xi,f_\xi,P_{s,\xi}\), and \(\psi_\xi\) are jointly
\(C^1\) in compatible stable-bundle coordinates, together with uniform
stable-coordinate domain containment, endpoint margins, and a bound
\(\partial_JH(\xi,J)\le-m_J<0\).  Under those additional hypotheses the
implicit-function theorem gives
\begin{equation}
 D_\xi J_c=-\frac{D_\xi H}{\partial_JH}.               \label{eq:Jc-derivative}
\end{equation}
None of these parameter-uniform graph and crossing hypotheses is supplied by
the center-parameter proposition.  Even that proposition proves a selected
stable-sheet intersection, not a biological onset.
To obtain the latter, invariant exit cones must carry the two local sides to
compact exit faces, and finite-time enclosures must attach those faces to the
quiet and outer attraction tubes.  The strict quiet capture at \(J=0.30\)
supplies one anchor; the corresponding outer-side tube is not yet complete.


